\section{Discussion and Limitations}\label{sec:discussion}
The central question is how to adapt VLA-based control when an existing action interface does not provide useful behavior on a new robot. The proposed sequence gives each learning phase a distinct role: dense reward helps discover a skill, binary reward defines its completion, self-imitation consolidates successful experience, and staged RL learns to complete the entire instruction. This separation also clarifies the measurements needed to assess the approach: skill acquisition, skill retention after consolidation, and final full-task success.

The cable-driven manipulator is an embodiment-adaptation test case. Its unusual mechanics make it useful for examining whether pretrained representations can support a newly learned controller. Success on this platform is evidence for that case, not proof of applicability to every manipulator. The complementary LIBERO study tests stage-aware feedback with an already trained action policy; it does not establish adaptation after head reinitialization. A stronger cross-embodiment claim requires applying the same training schedule to additional robot interfaces under declared evaluation protocols.

The evaluated SmolVLA system includes a learned residual action module and retained VLA components. Reinitializing an entire pretrained action expert is a distinct intervention and must be documented if claimed. Likewise, skill-level self-imitation should not be interpreted as a guarantee that another supervised pass on full-task successes will improve an already competent controller: the retained campaign includes unsuccessful full-task SFT attempts. Retention depends on the data mixture, task coverage, and action consistency.

The current simulation endpoint uses development scenes and one training lineage. Matched schedule, self-imitation, and composition ablations, independent seeds, and a locked final test remain necessary for stronger attribution. The LIBERO draft lacks trial denominators and quantitative throughput records. Physical success remains the author's drafting assumption until completed trial evidence is inserted. Manual rewards, historical scripted initialization, calibration, and inherited training costs also remain part of the method's scope.
